% GENERATED FILE. Edit canonical lessons, then run npm run generate:books.
% canonical-source-hash: fnv1a64:1df759407c0813fd
% canonical-lessons: SA-C272-offer-condolence, SA-C272-condolence-first-retrieval, SA-C272-condolence-final-retrieval

\chapter{One condolence, held as a whole}
\label{ch:sa-bounded-condolence}
\hlchaptermodality{272}

\begin{chapteropening}
\textbf{By the end of this chapter:} \emph{I can offer one attested condolence to bereaved family members.}

Complete chapter 272 by independently retrieving the bounded condolence a second time.
\end{chapteropening}

\section[\texorpdfstring{\sk{मम} \sk{संवेदनाः} … (mama saṃvedanāḥ …)}{मम संवेदनाः … (mama saṃvedanāḥ …)}]{One condolence, held as a whole}

\label{lesson:SA-C272-offer-condolence}



\begin{quote}
Recall \textbf{\sk{मम}}, “my”; \textbf{\sk{दुःखम्}}, “sorrow”; \textbf{\sk{शोकः}}, “grief”; \textbf{\sk{म्रियते}}, “dies”; and \textbf{\sk{दया}}, “compassion.” These known words prepare the meaning. They are not yet a condolence you can say to somebody.
\end{quote}

\subsection*{You'll want to know: one attested condolence}
\begin{quote}
\textbf{\sk{मम} \sk{संवेदनाः} \sk{शोकसन्तप्त}-\sk{परिजनैः} \sk{सह} \sk{सन्ति।}}
\end{quote}

\begin{quote}
\emph{mama saṃvedanāḥ śokasantapta-parijanaiḥ saha santi.}
\end{quote}

\begin{quote}
“My condolences are with the bereaved family members.”
\end{quote}

A contemporary Sanskrit report from \href{https://sanskrit.hindusthansamachar.in/Encyc/2026/5/28/Chief-Minister-Expresses-Grief-Demise-Poet-Dr-Bash.php}{\emph{Hindusthān Samācār}} quotes the same complete sentence, including the editorial hyphen inside \textbf{\sk{शोकसन्तप्त}-\sk{परिजनैः}}. The report types an ASCII colon after \textbf{\sk{परिजनै}}; the lesson normalizes that one character to the Devanagari visarga \textbf{\sk{ः}} shown above. All India Radio Sanskrit bulletins independently use plural \textbf{\sk{शोकसंवेदनाः}} for condolences directed to bereaved families. The sentence here is therefore attested contemporary Sanskrit, not a translation silently built from an English formula.

\begin{grammarlens}[title={Grammar Lens: keep the formula whole}]
Hear the sentence in three meaning-chunks:

\begin{itemize}
\raggedright
  \item \textbf{\sk{मम} \sk{संवेदनाः}} — “my condolences”
  \item \textbf{\sk{शोकसन्तप्त}-\sk{परिजनैः} \sk{सह}} — “with the bereaved family members”
  \item \textbf{\sk{सन्ति}} — “are”
\end{itemize}

For now, keep these chunks together as one whole social formula. This lesson does not teach the instrumental plural or license new sentences from its pieces.

The linked report also says \textbf{\sk{तस्मै} \sk{विनम्र}-\sk{श्रद्धाञ्जलिः}}, a respectful tribute directed to the person who died. That is a different act. The sentence you are learning offers a condolence to the bereaved family; it is not a tribute to the person who died.
\end{grammarlens}

\subsection*{Your turn}
A family has lost somebody close. Address the bereaved family members and say the complete condolence slowly once.

\begin{tcolorbox}[breakable,colback=teal!4,colframe=teal!35!black,title={Before you move on}]
Say, “My condolences are with the bereaved family members.” Keep the attested sentence whole.
\end{tcolorbox}
\section[(independent condolence retrieval …]{Offer the condolence: round one}

\label{lesson:SA-C272-condolence-first-retrieval}



\begin{quote}
Prepare to retrieve the complete sentence from its social purpose.
\end{quote}

\subsection*{Your turn: independent retrieval}
No Sanskrit model is shown before you answer.

A neighbour has lost a close relative. You are speaking to the bereaved family members, not paying tribute to the person who died. Say the Sanskrit for “My condolences are with the bereaved family members.”

\begin{grammarlens}[title={What you've built: check after answering}]
Now check: \textbf{\sk{मम} \sk{संवेदनाः} \sk{शोकसन्तप्त}-\sk{परिजनैः} \sk{सह} \sk{सन्ति।}} (\emph{mama saṃvedanāḥ śokasantapta-parijanaiḥ saha santi.})
\end{grammarlens}

\begin{tcolorbox}[breakable,colback=teal!4,colframe=teal!35!black,title={Before you move on}]
You have retrieved one bounded condolence once without a visible model. This is not a claim of a wider condolence repertoire or general social competence.
\end{tcolorbox}
\section[(independent condolence retrieval …]{Offer the condolence: round two}

\label{lesson:SA-C272-condolence-final-retrieval}



\begin{quote}
Retrieve the same social formula in a different setting.
\end{quote}

\subsection*{Your turn: independent retrieval in a fresh context}
No Sanskrit model is shown before you answer.

At a memorial gathering, you turn from remembering the person who died to address the bereaved family members directly. Say the Sanskrit for “My condolences are with the bereaved family members.”

\begin{grammarlens}[title={What you've built: check after answering}]
Now check: \textbf{\sk{मम} \sk{संवेदनाः} \sk{शोकसन्तप्त}-\sk{परिजनैः} \sk{सह} \sk{सन्ति।}} (\emph{mama saṃvedanāḥ śokasantapta-parijanaiḥ saha santi.})
\end{grammarlens}

\begin{tcolorbox}[breakable,colback=teal!4,colframe=teal!35!black,title={Before you move on}]
You can retrieve this one attested condolence twice without a visible model. That is not evidence of general social competence, human validation, learner attainment, or exam readiness.
\end{tcolorbox}
